\section{The horizon and all enclosing surfaces}\label{sec:geometry}

The estimates for the conformal factor make the coordinate spheres useful
barriers.  Their positive expansion outside the boundary also gives a closed
two-form that compares the boundary area with every enclosing cut.

\begin{proposition}\label{prop:geometry}
For all sufficiently large $L$, uniformly for $0\le s\le1$, the data constructed
in \Cref{sec:construction,sec:elliptic} are smooth source-free electrovacuum
data on
\[
 \Omega=\{x\in\R^3:r\ge1\},\qquad S=\{r=1\},
\]
with $\tr_gK=0$.  The metric is complete with its boundary included, has one
asymptotically flat end, and satisfies
\begin{equation}\label{eq:geometric-falloff}
 g_{ij}-\delta_{ij}=O_4(r^{-1}),\qquad K_{ij}=O_3(r^{-3}),\qquad
 \mathcal E^i,\mathcal B^i=O_3(r^{-2}).
\end{equation}
The standard effective $U(1)$ action preserves all the data and $S$, and
$\mu$ and $|J_{\mathrm{con}}|_g$ are integrable.  With the normal toward the end,
$S$ is a future MOTS, and no compact smooth embedded enclosing surface
in $\operatorname{int}\Omega$, possibly disconnected, has everywhere
nonpositive future outward expansion.  Every enclosing cut in the sense of
\Cref{sec:introduction} has area at least $A=\Area_g(S)$.  Moreover, for a
constant $C_A$ independent of $s,L$,
\begin{equation}\label{eq:area-estimate}
 64\pi\le A\le64\pi+C_A\frac{s^2}{L^2},
 \qquad A=64\pi+O(s^2/L^2),
\end{equation}
and $A(s,L)$ is continuous in $s$ for fixed $L$.
\end{proposition}

\begin{proof}
The domain is connected and oriented by the standard orientation of $\R^3$;
its boundary is a connected two-sphere.  Outside each sufficiently large
ball it is a single product end.  Smoothness up to $S$ and across the rotation
axis follows from \Cref{lem:seeds,prop:conformal-factor}; the constraints and
maximality follow from \Cref{lem:conformal}.  The seeds are invariant under
the standard rotations, and uniqueness of the fixed point makes $u$ invariant
as well.  Thus these rotations preserve $g,K,\mathcal E,\mathcal B$, with
generator $\eta=\partial_\phi$ of period $2\pi$.  The action is effective
because a rotation fixing every point of $\Omega$ has angle $0$ modulo $2\pi$.

The bound $1\le u\le M$ gives
$\delta\le g\le M^4\delta$.  The Euclidean intrinsic length metric on the
closed exterior is complete: a Cauchy sequence has a Euclidean limit in
$\Omega$, and local smooth boundary charts show convergence also in the
intrinsic length metric.  Comparison of curve lengths now proves completeness
for $g$ with $S$ included.  The derivative bounds in
\Cref{lem:seeds,prop:conformal-factor} and the product rule give
\eqref{eq:geometric-falloff}.

The electrovacuum constraints identify
\[
 \mu=\frac{|\mathcal E|_g^2+|\mathcal B|_g^2}{8\pi}=O(r^{-4}),
 \qquad
 J_{\mathrm{con}}=\frac{(\mathcal E\times_g\mathcal B)^{\flat_g}}{4\pi}.
\]
The fields vanish for $r\le L$ and are radial, hence parallel, for $r\ge2L$.
Consequently $J_{\mathrm{con}}$ is supported in the compact annulus
$L\le r\le2L$.  Since $\dd V_g=u^6\dd x\le M^6\dd x$, both required
densities are integrable.  The horizon and enclosing-surface assertions are
proved in \Cref{lem:positive-spheres,lem:area-minimum} below.
\end{proof}

\begin{lemma}\label{lem:positive-spheres}
For $L$ sufficiently large, the coordinate sphere $S_r=\{r=\mathrm{const}\}$
has future outward expansion zero at $r=1$ and strictly positive at every
$r>1$.  No compact smooth embedded enclosing surface in
$\operatorname{int}\Omega$, with any finite number of components, is weakly
future outer trapped.
\end{lemma}

\begin{proof}
The normal to $S_r$ toward infinity is $\nu_r=u^{-2}\partial_r$.  Since
$\sigma$ has only mixed radial--azimuthal components,
$K|_{TS_r\times TS_r}=0$, including on the axis by smoothness.  The mean
curvature is therefore also the future outward expansion, and direct
divergence gives
\begin{equation}\label{eq:sphere-expansion}
 \theta_+(S_r)=H_{S_r}
 =\frac{1}{u^6r^2}\partial_r(r^2u^4)
 =\frac{2W}{ru^3},\qquad W=u+2ru_r.
\end{equation}
This formula permits angular dependence of $u$: the normal has no angular
components.  The Robin condition \eqref{eq:robin} gives $W(1,\omega)=0$.
For $r\ge2$, \eqref{eq:radial-bounds} gives
\[
 W=1-r^{-1}+v+2rv_r\ge\frac12-\frac{Cs^2}{L^2}.
\]
On $1\le r\le2$, the same estimates yield
\[
 W_r=r^{-2}+3v_r+2rv_{rr}\ge\frac14-\frac{Cs^2}{L^2}.
\]
Increasing $L$ once, uniformly over $s\in[0,1]$, we obtain
\begin{equation}\label{eq:positive-W}
 \begin{aligned}
 W_r\ge\frac18,\qquad W\ge\frac{r-1}{8}
     &\quad(1\le r\le2),\\
 W\ge\frac14&\quad(r\ge2).
 \end{aligned}
\end{equation}
Thus the strict positivity holds even arbitrarily close to $S$.

Let $\Sigma\subset\operatorname{int}\Omega$ be any compact smooth embedded
enclosing surface, without a symmetry or connectedness restriction.  To
cover also the interpretation allowing components with interior edges, let
$\Sigma_{\mathrm{cl}}$ be the union of its components without boundary.
The other components do not affect separation.  Indeed, a path avoiding
$\Sigma_{\mathrm{cl}}$ can be perturbed to miss their edges and cross their
interiors transversely at finitely many points.  For each crossing, choose
an embedded arc in that component from the crossing to an edge.  A narrow
tubular strip along the arc gives a detour on one side of the component,
around the edge, and back on the other side, removing the crossing without
introducing another.  Compactness and disjointness of the components, and
their location in $\operatorname{int}\Omega$, let every detour avoid the
other components and $S$.  Thus any path avoiding $\Sigma_{\mathrm{cl}}$
can avoid all of $\Sigma$, so $\Sigma_{\mathrm{cl}}$ still encloses $S$
and is nonempty.

Choose $p\in\Sigma_{\mathrm{cl}}$ where $r$ has its global maximum
$R>1$ on $\Sigma_{\mathrm{cl}}$.  The increasing radial ray beyond $p$
misses $\Sigma_{\mathrm{cl}}$.  Starting just beyond $p$, the same detours
past any edged components connect this side to infinity without meeting
$\Sigma$.  Hence the prescribed normal at $p$ is
$\nu_\Sigma(p)=u(p)^{-2}\partial_r$, and $T_p\Sigma=T_pS_R$.
With the convention $H_\Sigma=\operatorname{div}_\Sigma
\nu_\Sigma$, restriction of the ambient Hessian gives, at $p$,
\[
 0\ge\Delta_\Sigma r
 =\tr_{T_p\Sigma}\nabla_g^2r-H_\Sigma\,\nu_\Sigma(r)
 =|\nabla_g r|_g\bigl(H_{S_R}-H_\Sigma\bigr).
\]
For the last identity, the tangential Hessian of a defining function for a
level surface equals its second fundamental form times the length of its
gradient.  It follows that $H_\Sigma(p)\ge H_{S_R}(p)>0$.  Tangency also gives
$\tr_\Sigma K(p)=0$, so $\theta_+(\Sigma)(p)>0$.  This excludes precisely
the condition $\theta_+\le0$ everywhere.
\end{proof}

\begin{lemma}\label{lem:area-minimum}
For the same choices of $L$, every enclosing cut $\Gamma=\partial D$ from the
specified class has $\Area_g(\Gamma)\ge\Area_g(S)$.  Thus $S$ attains the
minimum enclosing area, and \eqref{eq:area-estimate} holds.
\end{lemma}

\begin{proof}
Let $\pi:\Omega\longrightarrow\Sph^2$ be the radial projection
$\pi(r,\omega)=\omega$, and let $\dd A_{\Sph^2}$ denote the unit round area
form.  Define the smooth closed two-form
\[
 \beta=\pi^*\bigl(u(1,\omega)^4\dd A_{\Sph^2}\bigr).
\]
It is closed because every two-form on $\Sph^2$ is closed.  Its comass, namely
the supremum of $|\beta(e_1,e_2)|$ over $g$-orthonormal pairs, is
\begin{equation}\label{eq:calibration-comass}
 \|\beta\|_{*,g}(r,\omega)
 =\frac{u(1,\omega)^4}{r^2u(r,\omega)^4}\le1.
\end{equation}
Indeed, $\beta$ annihilates radial vectors and its value on an oriented
orthonormal angular pair is the displayed ratio.  Furthermore,
\[
 \partial_r(ru^2)=u(u+2ru_r)=uW>0\qquad(r>1),
\]
by \Cref{lem:positive-spheres}; hence $ru(r,\omega)^2\ge u(1,\omega)^2$.
The comass estimate controls every tangent plane, so it does not require
$\Gamma$ to be a radial graph.

For completeness, take any connected smooth codimension-zero submanifold
$D\subset\Omega$ that is closed in $\Omega$, has interior contained in
$\operatorname{int}\Omega$, contains the entire sufficiently distant part
of the end, and has compact smooth embedded two-sided intrinsic boundary
$\Gamma$.  Choose $R$ beyond that boundary and beyond the radius where $D$
contains the entire end.  Then
\[
 D_R=D\cap\{r\le R\}
 \quad\text{is compact and smooth, with}\quad
 \partial D_R=\Gamma\sqcup S_R.
\]
Near $\Gamma$ the truncation leaves $D$ unchanged, including at contacts
with $S$; near $S_R$ it is a regular level-set truncation in the interior
of $D$.  These boundary pieces are disjoint, so $D_R$ has no corners.
Here $D_R$ inherits the orientation of $\Omega$, and its boundary is the
\emph{full intrinsic boundary}.  In particular, any portion coinciding with
$S$ already belongs to $\Gamma$ and contributes to this boundary integral.
Stokes' theorem with the induced boundary orientations gives
\[
 \int_\Gamma\beta=-\int_{S_R}\beta
 =-\int_{\Sph^2}u(1,\omega)^4\dd A_{\Sph^2}=-A.
\]
Using \eqref{eq:calibration-comass}, including all components, we conclude
\[
 A=\left|\int_\Gamma\beta\right|
 \le\int_\Gamma\|\beta\|_{*,g}\,\dd A_g
 \le\Area_g(\Gamma).
\]
The choice $D=\Omega$ is allowed and gives $\Gamma=S$, so the infimum equals
$A$ and is attained.  No symmetry or finite-perimeter approximation has
entered this comparison.

Finally, $u(1,\omega)=2+v(1,\omega)$ with
$0\le v(1,\omega)\le Cs^2/L^2$.  Since $u\le M$,
\[
 0\le A-64\pi
 =\int_{\Sph^2}\bigl((2+v)^4-16\bigr)\dd A_{\Sph^2}
 \le4M^3\int_{\Sph^2}v\,\dd A_{\Sph^2}
 \le C_A\frac{s^2}{L^2}.
\]
This proves \eqref{eq:area-estimate}.  Supremum-norm continuity of $v$ in $s$
from \Cref{prop:conformal-factor} also proves continuity of $A(s,L)$.
\end{proof}
